\documentclass[12pt,a4paper]{article}
\usepackage[utf8]{inputenc}
\usepackage[T1]{fontenc}
\usepackage{geometry}
\usepackage{setspace}
\usepackage{booktabs}
\usepackage{longtable}
\usepackage{array}
\usepackage{graphicx}
\usepackage{hyperref}
\usepackage{enumitem}
\usepackage{titlesec}
\usepackage{xcolor}

\geometry{margin=1in}
\onehalfspacing
\setlist{nosep}
\hypersetup{
  colorlinks=true,
  linkcolor=blue,
  urlcolor=blue,
  citecolor=blue
}

\titleformat{\section}{\large\bfseries}{\thesection}{0.6em}{}
\titleformat{\subsection}{\normalsize\bfseries}{\thesubsection}{0.6em}{}

\begin{document}

\begin{titlepage}
    \centering
    {\Large Concordia University\par}
    \vspace{0.3cm}
    {\large Department of Computer Science and Software Engineering\par}
    \vspace{1.4cm}
    {\Huge \textbf{SOEN 345 Project Report}\par}
    \vspace{0.4cm}
    {\Large Cloud-based Ticket Reservation Application\par}
    \vspace{1.4cm}

    \begin{tabular}{>{\bfseries}p{4.4cm}p{8.2cm}}
        Student Name: & Harjot Minhas \\
        Student ID: & 40315397 \\
        Course: & SOEN 345 -- Software Testing, Verification and Quality Assurance \\
        Term: & Winter 2026 \\
        Submission Date: & April 13, 2026 \\
        Public GitHub Repository: & \url{https://github.com/M1nha5/soen345-ticket-reservation} \\
    \end{tabular}

    \vfill
    {\large Final Submission\par}
\end{titlepage}

\tableofcontents
\newpage

\section{Project Overview}

This project implements a cloud-based Android ticket reservation application where users can browse events, reserve tickets, cancel reservations, and receive confirmation records. The system supports multiple roles:

\begin{itemize}
    \item Customer
    \item Organizer
    \item Admin
\end{itemize}

The implementation follows the course requirement to use Java and demonstrates software engineering, testing, and quality assurance practices with a Firebase cloud backend.

\subsection{Project Objectives}

\begin{itemize}
    \item Build a mobile ticket reservation solution using Java on Android.
    \item Implement role-based workflows for customer, organizer, and admin.
    \item Provide cloud-based storage and authentication.
    \item Validate correctness through unit, component, functional, and acceptance testing.
    \item Use professional engineering tools: GitHub, CI/CD, Android Studio, JUnit, and Firebase.
\end{itemize}

\section{Software Development Method}

\subsection{Adopted Methodology}

The project used an iterative Scrum-style development process. Work was delivered in short increments with frequent integration and testing.

\subsection{Lifecycle and Iterations}

\textbf{Iteration 1: Foundation}
\begin{itemize}
    \item Android project setup
    \item Firebase integration (Auth, Firestore, Functions)
    \item Core data models and service layer
\end{itemize}

\textbf{Iteration 2: Core Functional Features}
\begin{itemize}
    \item Email/password and phone OTP authentication
    \item Event listing, filtering, reservation and cancellation
    \item Organizer registration and admin approval flow
    \item Event management (add/edit/cancel)
\end{itemize}

\textbf{Iteration 3: Hardening and Quality}
\begin{itemize}
    \item Firestore rules hardening
    \item Confirmation function reliability improvements
    \item UI improvements and empty-state handling
    \item CI stability improvements and regression testing
    \item Final release build signing and deployment readiness
\end{itemize}

\subsection{Definition of Done}

A work item was considered done only when:
\begin{itemize}
    \item implementation was complete in app/backend,
    \item feature behavior was validated by tests or acceptance checks,
    \item integration worked with Firebase,
    \item documentation was updated in \texttt{docs/}.
\end{itemize}

\section{Requirements Analysis}

\subsection{Functional Requirements Coverage}

\begin{longtable}{p{0.7cm}p{7.2cm}p{2.5cm}p{4.2cm}}
\toprule
\textbf{ID} & \textbf{Requirement} & \textbf{Status} & \textbf{Implementation Evidence} \\
\midrule
FR1 & User registration with email or phone number & Implemented & Email registration and phone OTP login screens \\
FR2 & View list of available events & Implemented & Event list screen and Firestore \texttt{events} collection \\
FR3 & Search/filter events by date, location, category & Implemented & In-app filter logic and UI controls \\
FR4 & Cancel reservations & Implemented & Reservation list with cancel action and transaction logic \\
FR5 & Receive confirmation via email/SMS channel record & Implemented & Callable Cloud Function writes confirmation outcome \\
FR6 & Admin/Organizer add new event & Implemented & Event editor and role-based access \\
FR7 & Admin/Organizer edit existing event & Implemented & Event editor update flow \\
FR8 & Admin/Organizer cancel event & Implemented & Event status transition to cancelled \\
\bottomrule
\end{longtable}

\subsection{Non-Functional Requirements Coverage}

\begin{longtable}{p{0.7cm}p{6.7cm}p{2.4cm}p{4.8cm}}
\toprule
\textbf{ID} & \textbf{Requirement} & \textbf{Status} & \textbf{Implementation Strategy} \\
\midrule
NFR1 & Support concurrent users without performance degradation & Implemented & Firestore transactions for reservation/cancellation preserve inventory consistency under concurrent access \\
NFR2 & Cloud-based high availability & Implemented & Firebase Auth, Firestore, and Cloud Functions managed cloud infrastructure \\
NFR3 & UI should be simple and user-friendly & Implemented & Role-based dashboard, improved visual hierarchy, clearer forms, and explicit empty states \\
\bottomrule
\end{longtable}

\section{System Architecture and Design}

\subsection{Architecture Overview}

The application follows a client-cloud architecture:
\begin{itemize}
    \item Android app (presentation + client logic)
    \item Firebase Authentication (identity)
    \item Cloud Firestore (data persistence)
    \item Firebase Cloud Functions (server-side confirmation processing)
\end{itemize}

\begin{figure}[h!]
\centering
\fbox{
\parbox{0.92\linewidth}{
\centering
\textbf{Architecture Diagram Placeholder}\\
\vspace{0.2cm}
Android App $\leftrightarrow$ Firebase Auth\\
Android App $\leftrightarrow$ Firestore\\
Android App $\leftrightarrow$ Cloud Functions\\
Cloud Functions $\leftrightarrow$ Firestore
}}
\caption{High-level system architecture.}
\end{figure}

\subsection{Use Case Design}

\textbf{Primary actors:}
\begin{itemize}
    \item Customer
    \item Organizer
    \item Admin
\end{itemize}

\textbf{Customer use cases:}
\begin{itemize}
    \item Register/Login
    \item Browse events
    \item Search/filter events
    \item Reserve tickets
    \item Cancel reservation
\end{itemize}

\textbf{Organizer/Admin use cases:}
\begin{itemize}
    \item Create event
    \item Edit event
    \item Cancel event
\end{itemize}

\textbf{Admin-only use cases:}
\begin{itemize}
    \item Approve/reject organizer accounts
\end{itemize}

\subsection{Database Design (Firestore)}

\textbf{Collection: \texttt{users}}
\begin{itemize}
    \item \texttt{uid}, \texttt{name}, \texttt{email}, \texttt{phone}, \texttt{role}, \texttt{status}, \texttt{createdAt}
\end{itemize}

\textbf{Collection: \texttt{events}}
\begin{itemize}
    \item \texttt{id}, \texttt{title}, \texttt{category}, \texttt{location}, \texttt{dateTimeMillis}
    \item \texttt{totalTickets}, \texttt{availableTickets}, \texttt{organizerId}, \texttt{status}, \texttt{createdAt}
\end{itemize}

\textbf{Collection: \texttt{reservations}}
\begin{itemize}
    \item \texttt{id}, \texttt{userId}, \texttt{eventId}, \texttt{eventTitle}, \texttt{eventLocation}
    \item \texttt{eventDateTimeMillis}, \texttt{tickets}, \texttt{status}, \texttt{createdAt}
\end{itemize}

\textbf{Collection: \texttt{confirmations}}
\begin{itemize}
    \item \texttt{reservationId}, \texttt{userId}, \texttt{eventId}, \texttt{action}
    \item \texttt{smsStatus}, \texttt{emailStatus}, \texttt{createdAt}
\end{itemize}

\subsection{Core Sequence (Reservation)}

\begin{enumerate}
    \item Customer selects an event and ticket quantity.
    \item Android app requests reservation through service layer.
    \item Firestore transaction loads latest event document.
    \item System validates availability and decrements available tickets atomically.
    \item Reservation document is created in the same transaction.
    \item App calls \texttt{sendConfirmation} Cloud Function.
    \item Cloud Function writes confirmation outcome to \texttt{confirmations}.
\end{enumerate}

\section{Implementation Summary}

\subsection{Role-Based Behavior}

\begin{itemize}
    \item \textbf{Customer:} browse/filter events, reserve/cancel tickets, view reservations.
    \item \textbf{Organizer:} create/edit/cancel events only after admin approval.
    \item \textbf{Admin:} approve/reject organizers and manage events.
\end{itemize}

\subsection{Security and Data Integrity}

\begin{itemize}
    \item Firestore rules enforce role-based authorization.
    \item Rules prevent self role/status privilege escalation.
    \item Reservation/event updates use constrained write rules and transaction checks.
    \item Cloud Function callable endpoint verifies authenticated caller and permissions.
\end{itemize}

\subsection{User Interface}

\begin{itemize}
    \item Redesigned, user-friendly screens for authentication, dashboards, and event lists.
    \item Clear empty states for no events and no reservations.
    \item Improved reservation input using number picker for ticket count.
    \item Phone input normalization for North American numbers (auto \texttt{+1} support).
\end{itemize}

\section{Software Testing Method and Results}

\subsection{Testing Strategy}

Testing was performed across four levels:
\begin{itemize}
    \item Unit tests for policy/filter logic
    \item Component tests for booking flow logic
    \item Functional UI checks using Espresso
    \item Manual acceptance tests for complete end-to-end scenarios
\end{itemize}

\subsection{Unit and Component Tests}

\textbf{Execution command:}
\begin{quote}
\texttt{./gradlew testDebugUnitTest}
\end{quote}

\textbf{Result summary (April 13, 2026):}
\begin{itemize}
    \item Total: 10
    \item Passed: 10
    \item Failed: 0
    \item Errors: 0
    \item Skipped: 0
\end{itemize}

\textbf{Covered test suites:}
\begin{itemize}
    \item Event filtering logic
    \item Role policy and organizer approval constraints
    \item Reservation policy validation
    \item Booking component flow consistency
\end{itemize}

\subsection{Functional and Acceptance Tests}

\textbf{Functional UI test:}
\begin{itemize}
    \item Espresso test verifies login screen actions and visibility.
\end{itemize}

\textbf{Acceptance scenarios executed:}
\begin{itemize}
    \item Customer registration and login (email/password)
    \item Organizer registration (pending status)
    \item Admin approval/rejection for organizers
    \item Event creation/update/cancellation by authorized role
    \item Customer browsing/filtering/reserving/cancelling
    \item Phone OTP authentication flow
    \item Confirmation generation on reservation state changes
\end{itemize}

\textbf{Artifacts:}
\begin{itemize}
    \item \texttt{docs/testing-plan.md}
    \item \texttt{docs/acceptance-results.md}
\end{itemize}

\section{Development Tools and Collaboration}

\begin{longtable}{p{4.2cm}p{10cm}}
\toprule
\textbf{Category} & \textbf{Tooling Used} \\
\midrule
IDE & Android Studio \\
Programming Language & Java (Android app), JavaScript (Cloud Functions) \\
Version Control & Git + GitHub \\
CI/CD & GitHub Actions (\texttt{.github/workflows/android-ci.yml}) \\
Testing Frameworks & JUnit, Mockito, Espresso \\
Cloud Platform & Firebase Authentication, Firestore, Cloud Functions \\
\bottomrule
\end{longtable}

The repository was continuously updated on the \texttt{main} branch with feature-based commits and verification passes after major changes.

\section{Rubric Mapping}

\begin{longtable}{p{5.8cm}p{2.5cm}p{7cm}}
\toprule
\textbf{Rubric Area} & \textbf{Weight} & \textbf{Coverage in This Report} \\
\midrule
Requirements Analysis \& System Design & 15 & Sections 4 and 5 (requirements tables, architecture/design, Firestore model, sequence flow) \\
Functional/Non-Functional Requirements & 15 & Sections 4 and 6 (implementation and NFR strategy) \\
Software Testing \& Quality Assurance & 40 & Section 7 (test strategy, results, acceptance evidence) \\
Professional Report Quality & 30 & Structured sections, formal language, tables, traceability, GitHub evidence \\
\bottomrule
\end{longtable}

\section{Conclusion}

The project objectives were achieved with a complete cloud-based ticket reservation system that satisfies the specified functional and non-functional requirements. The final system includes role-based workflows, secure backend integration, transactional reservation consistency, and test-backed quality evidence. Continuous integration and comprehensive documentation support maintainability and reproducibility.

\section*{Appendix A: GitHub and Build Information}

\begin{itemize}
    \item Public repository: \url{https://github.com/M1nha5/soen345-ticket-reservation}
    \item Main Android package: \texttt{com.harjot.ticketreservation}
    \item Firebase project id: \texttt{soen345project-44582}
    \item Release artifact path: \texttt{app/build/outputs/apk/release/app-release.apk}
\end{itemize}

\section*{Appendix B: Suggested Screenshot List for Final PDF}

\begin{enumerate}
    \item Login and registration screens
    \item Phone OTP authentication screen
    \item Customer event list with filters
    \item Reservation creation flow
    \item Reservation cancellation result
    \item Organizer event management page
    \item Admin organizer approval page
    \item Firestore collections snapshot (\texttt{users/events/reservations})
    \item Confirmation records snapshot
    \item GitHub Actions successful build/test run
\end{enumerate}

\end{document}

